% !TeX root = ../papers/nc-p3.tex
\section{直線と横断二次曲線の三交換子と一次元ファイバー}
\label{sec:marked-transverse-conic-commutators}
\begin{lemma}\label{lem:marked-squarefree-symmetrizer}
$\det A$ が被約な平面三次曲線なら $\Gamma_Q$ が $A$ の行の消去空間に入る
線型解空間は次元 $6+r$ である。$r$ は行空間内の射影的な階数一テンソルの個数であり、
各々は行列式の相異なる直線成分に対応する。実際の三交換子表示で $\det C_Q\ne0$ なら、
そのようなテンソルはちょうど、$Q_n$ がスカラーと階数一行列の和になる方向 $[n]$ に一つずつ対応する。
\end{lemma}
\begin{proof}
定理\ref{thm:marked-elliptic-commutator-doubling}の積写像の双対で、一つの評価を恒等、
残りを $B_1,B_2$ とする。一般直線は被約三次曲線と相異なる三点で横断するので、
一般の $B_1,B_2$ の線型結合は固有値が相異なる。その固有基底で三因子間の作用の移動ができる
対称三テンソルは $\sum c_i e_i^{\otimes3}$ に限る。
両作用とも移動できることから $c_i\ne0$ なら $e_i$ は共通固有ベクトルであり、
逆にそのような立方は全て余核の双対に入る。各共通固有ベクトルは評価を恒等にした
行空間の階数一テンソルを与え、共通固有文字は一般結合上で既に相異なる。
最後に $n\otimes b$ が行空間に入る条件は全ての $h\in U^*$ に対し
$b,h,Q_n^t h$ が従属することである。$h$ と $h+b$ を比較すると $Q_n^t b$ は $kb$ に入る。
商 $U^*/kb$ の全直線を保つので商上でスカラーとなり、$Q_n$ は所要の形である。
スカラーだけなら二つの交換子が消えるから除かれる。階数一補正は一意である。
\end{proof}

\begin{theorem}\label{thm:marked-transverse-conic-commutators}
標数零の代数閉体で $\det C$ が直線とそれに横断する滑らかな二次曲線の積であるとする。
直線を $x=0$ とする座標を選ぶ。順序付き三交換子の全逆像は次の互いに素な二部分から得られる。
\begin{enumerate}
\item $Q_0$ に階数一補正がない部分が存在するのは、目標が両交点で局所自由で
右核の成分次数が $(1,2)$ のときに限り、その条件は十分でもある。
逆像のこの部分は四点、二点、または二点と一つのアフィン二次曲面である。
\item 階数一補正を持つ部分は目標の直線上の鉛筆から高々二つの旗 $(K,kw)$ と以下の
四変数の行列階数判定、線型方程式、スカラー平方根で例外値を含めて得られる。
ここで $K$ は平面で、$w\in K$ は要求しない。両核が直線上で動き、旗が半単純な場合には、
固定倍率ごとにアフィン線型空間上の一つの二次式、または二行三列の階数一以下の条件で尽きる。
\end{enumerate}
局所自由で次数 $(1,2)$ の目標の一般点で、全逆像は四つの孤立解と
$\mathbf G_m$ に同型な二つのアフィン二次曲線である。
\end{theorem}
\begin{proof}[補助表示の被覆]
生成添字の基底変更で $E:x(yz-x^2)=0$ とし、成分を
$(0:s:t)$、$(st:s^2:t^2)$ で正規化する。前節と同様に
$T(-H)$ はEuler標数と $h^0$ が零の階数一捩れなし層である。
局所自由でない交点を分離した部分正規化からの線束の直像となる。
両方で自由なら $T$ の成分次数は $(0,3),(1,2),(2,1)$ のいずれかで、
両端の貼合せの倍率比 $\lambda\ne0$ が残る。$(1,2)$ では $\lambda=1$ を除く。
実際 $T(-H)$ の成分次数和は零、片方が2以上なら他方で零にして両交点で消える切断がある。
一次数 $1,-1$ のときそのような切断はなく、次数 $0,0$ では両定数の貼合せが判定する。
一交点だけ分離すれば次数和が一つ下がり $T$ は $(0,2),(1,1)$、貼合せ倍率は基底変更で消える。
二点を分離すれば $T$ の両次数は $(0,1)$ である。これで全ての可能性を覆う。

直線上の右・左核点の次数の和は、分離した交点の個数を $j$ として $2-j$ である。
余因子の二次式は両核の積で、階数一の交点では直線上にちょうど一次の共通零を持つ。
実際特異点の完備局所環で非自由な余核は二枝の正規化の直和であり、
その平面上の最小自由表示は二局所方程式の対角、残る単位ブロックは核に寄与しない。
従って左核が直線上で定点なら前補題の階数一の場合に入る。
左核が動き右核が定点の場合は、自由な $(0,3)$ または一交点非自由の $(0,2)$ である。
貼合せた切断 $H^0(T)$ の基底と双一次関係を取ると、これらの代表は
\[
 \begin{pmatrix}-x&y&0\\z&-x&0\\\lambda y&z&-x\end{pmatrix},
 \qquad \begin{pmatrix}-x&y&0\\0&-x&y\\0&z&-x\end{pmatrix}
\]
（または交点を交換）である。実際右核はそれぞれ直線で定点、二次曲線で
$(s^2t:st^2:\lambda s^3+t^3)$、$(s^2:st:t^2)$ を与える。
第二表示では直線の定点は $e_0$ であり、二次曲線の二端のうち一端だけこれに一致する。
両表示で六線型解を解くと全 $Q_i$ がそれぞれ $ke_2,ke_0$ を保つ。
従って交換子鉛筆は恒等的に特異である。

残る自由な $(1,2)$ の代表とその交換子の参照行列は
\[
 A_\lambda=\begin{pmatrix}-x&y&0\\z&x(1/\lambda-1)&-y/\lambda\\0&z&-x\end{pmatrix},
 \quad M_\lambda=\begin{pmatrix}-\lambda x&y&0\\z&x(\lambda^2-1)/\lambda&-y\\0&-\lambda z&x\end{pmatrix}.
\]
$A_\lambda$ の核は直線で $(s:0:\lambda t)$、二次曲線で $(s^2:st:t^2)$
（直線上では第二行を用いる）。$M_\lambda$ の核は順に
$(s:0:t)$、$(s^2:\lambda st:\lambda^2t^2)$ であり、目標の貼合せは $\lambda^{-2}$ となる。
前節と同じ $Z$ の座標を使い、$Q$ の六成分
$((Q_1)_{12},(Q_2)_{01},(Q_2)_{02},-(Q_0)_{21}/\lambda,(Q_2)_{12},(Q_2)_{21})$ を
$(-U,-(\lambda-1)D-P,-B,A,D,(\lambda-1)T)$ とする。
方程式 $\operatorname{im}\Gamma_Q\subset W$ のこの解は $Z$ からの線型同型になる。
$C_Q(e_0)M_\lambda(e_0)^{-1}$ から前節と同じ六単項式の係数で作る
$\Psi_\lambda$ に対して、再び $C_Q=\Psi_\lambda(\operatorname{adj}Z)M_\lambda$。
行優先の六行 $0,1,2,3,4,6$ の小行列式は $-\lambda^4$ である。
六列は多項式であり、$\lambda=-1$ では $\det M_{-1}=0$ なので解がない。
その他は各平方根ごとに先の $C=L M R$ の手順が全ての実際の解を与える。
右核の二次曲線の全点を固定する射影変換は恒等なので共役の一意性も成り立つ。
他方の平方根の $S'$ は $(P,T,U,-A,-B,2P-D)$、
\[
      4\det\Psi_\lambda(S)=\lambda(\lambda+1)^2\det S-
                                   \lambda(\lambda-1)^2\det S' .
\]
少なくとも一方は可逆である。$S=E_{00}$ なら $\det S'=-4$ なので
階数一とその近傍の階数二も起こる。
以上の六空間・余因子恒等式・小行列式は
\path{computations/verify_transverse_conic_commutators.py} で検算できる。
\end{proof}

\begin{proof}[階数一補正の全復元]
$N=w b^t$、$K=\ker b^t$ とする。$[N,-]$ の像は
$\{Y:Y K\subset kw,\ \operatorname{tr}Y=0\}$ である。
包含は積から明らかで、両空間は四次元である（半単純型も冪零型も中心化環は
全行列内で五次元）。従って候補 $(K,kw)$ は $C_1K+C_2K\subset kw$ の全解である。

具体的には直線上の一般点の右核は常に $K$ に含まれ、左核は $w$ を消す。
各核の線型包を、余因子の共通因子を除いた多項式ベクトルの係数から求める。
どちらかが平面全体を射影的に張れば候補はない。右核のベクトル包が二次元なら
それが $K$ であり、$C_1K+C_2K$ が一次元のときだけその像を $kw$ とする。
左核の包が二次元のときは転置で同様である。両核が定点 $kr,k\ell$ なら
$U/kr\to\ker\ell^t$ の二行二列の正則鉛筆を得る。可逆な一スライスで割り、他の
スライスの各固有直線を $K/kr$ とすればよい。二スライスが比例なら
$\det C$ の直線に沿う一次係数は一次式の平方の倍になり横断性に反する。
これで旗は高々二つである。

旗ごとに $N$ の代表を一つ固定し、$Q_0=\rho(N-\operatorname{tr}(N)I/3)$、
$X_i=\rho Q_i$ $(i=1,2)$ とする。線型に $[N,H_1]=C_2$、$[N,H_2]=-C_1$ を解き、
跡零中心化環の基底を $Z_j$ $(1\le j\le4)$ とする。
四変数で $X_1=H_1+\sum t_j Z_j$ を自由に動かし、八行五列の拡大前の係数行列と右辺
\[
       D(t)=\big([X_1,Z_1],\ldots,[X_1,Z_4],-C_0\big),
              \qquad k(t)=-[X_1,H_2]
\]
を作る（成分を跡零八座標に並べる）。
$\operatorname{rank}D=\operatorname{rank}(D\mid k)$ を満たす階数層で
$D(v_1,\ldots,v_4,h)^t=k$ を線型に解き、$h\ne0$ の解を全て取る。
$X_2=H_2+\sum v_j Z_j$、$\rho=\pm\sqrt h$ で全逆像となる。
各層の非零小行列式を選ぶと自由変数と有理式で陽に書ける。
入力の任意の特殊化についても四変数だけの小行列式による条件であり、
$h\ne0$ を満たす解の有無は特解と線型核の最終座標で判定できる。
必要十分性は三つの交換子への代入による。

$w\notin K$ なら $N=\operatorname{diag}(1,0,0)$ とする。
$X_i$ を $1+2$ ブロックに分けて上右行を $\alpha_i$、下左列を $u_i$ とすれば
これらは最初の二つの既知の交換子で固定されている。
右下ブロックと左上スカラー倍の差を $D_i$ とすると、跡零条件により
任意の $D_i$ は対角部分を唯一に定める。$h=\rho^2$ は
$h(C_0)_{00}=\alpha_1u_2-\alpha_2u_1$ を満たす非零数を全て取る。残る方程式は
\[
 \alpha_1D_2-\alpha_2D_1=h(C_0)_{01},\quad
 D_1u_2-D_2u_1=h(C_0)_{10},\quad
 [D_1,D_2]=J=h(C_0)_{11}-u_1\alpha_2+u_2\alpha_1 ,
\]
ここで $00,01,10,11$ はブロックを表す。$J$ は跡零である。
$J\ne0$ なら $\operatorname{tr}(D_iJ)=0$ の二つを最初の線型条件に加える。
$\mathfrak{sl}_2$ の跡対合とブラケットが与える交代三形式は体積形式であり、
$J^\perp$ の二元の交換子は $kJ$ に入る。よって最後は $J$ の非零成分一つに
対する二次式だけである。$J=0$ なら $D_1,D_2$ の跡零部分が従属する条件となる。

最後に次数 $(1,2)$ の局所自由目標の空間は既約である。
固定した成分対で貼合せ倍率、右核座標の射影基底、跡零左倍率の六空間の可逆開部分を
取ればよく、横断する成分対も一つの射影軌道をなす。四つの孤立解の条件は既に稠密開である。
半単純な一旗、非零な $h$ の一意性、線型解の次元2と滑らかな二次曲線、相異なる二無限遠点は
いずれも以下の点で確認できる開条件である。$Q_0=\operatorname{diag}(2,-1,-1)/3$ と
\[
 Q_1=\begin{pmatrix}4&1&1\\-1&-2&3\\0&3&-2\end{pmatrix},\qquad
 Q_2=\begin{pmatrix}-1/3&-2&3\\1&2/3&-1\\-2&-3&-1/3\end{pmatrix}
\]
を取ると、目標の二次因子は $38x^2-260xy+308xz+27y^2-30yz-3z^2$、
$h=1$、$J=\left(\begin{smallmatrix}-6&-3\\3&6\end{smallmatrix}\right)$。
$D_2$ の下行を $(g,k)$ として線型消去すれば残る式は
$27g^2-24gk+234g-6k^2-140k+459=0$ になる。両核は直線で次数1であり、全交点で自由である
（$\alpha_i$ 同士と $u_i$ 同士が各々独立）。
従ってこの非空開条件を同時に満たす一般目標では、$\rho$ の二符号ごとに
射影二次曲線から二点を除いたものが得られる。
\end{proof}
